\begin{abstract}
Self-evolving language-model agents distill interaction trajectories into a
persistent memory bank, and the quality of that self-improvement loop depends on
how credit is assigned to individual memories. We argue that the credit signals
used by current systems are \emph{correlational}: a memory that merely co-occurs
with success is reinforced, so \emph{retrieval is mistaken for contribution}.
We hypothesize that under such credit \emph{superstitious memories}---entries with
zero or negative causal effect that nonetheless accumulate utility and occupy
retrieval slots---are not corrected, and that their fraction grows with
deployment length. The difficulty is that a memory's \emph{immediate conditional
retrieval contribution}---how much the immediate task outcome degrades when the
memory is removed from the retrieved set---is defined only counterfactually, and
per-memory intervention appears unaffordable in an online, non-stationary, closed
loop. We propose \method{},
which is designed to make this counterfactual
credit affordable through a strictly linear pipeline: Randomized Interventional
Trials (\rit{}) collect small-budget causal labels, Amortized Counterfactual
Attribution (\aca{}) compresses them into a per-event score that is constant in
bank size (at fixed feature dimension $d$ and top-$k$), and two
minimal consumers---threshold governance and scope-gated retrieval---read the
same signal. We plan a systematic evaluation testing whether \aca{} recovers
this immediate credit at low cost, whether its consumers flatten superstition and
reduce cross-task interference, and whether \method{} improves average success
over the strongest baseline; the planned headline gain is
\tbd{MAIN_AVG_GAIN_OVER_BEST_BASELINE}, pending verification.
\end{abstract}
